% Môn: Vật lý
% Lớp: 12
% Chương: Chương IV. Vật lí hạt nhân
% Bài: L12C4 Bài 22. Phản ứng hạt nhân và năng lượng liên kết · Tính toán nhiệt lượng vật thu vào hoặc tỏa ra
% ===== Câu 98 =====
% ID: L12C4B22-82-TL
% Mức: NB
\dangbt{Tính toán nhiệt lượng vật thu vào hoặc tỏa ra}
\begin{bt}
% ID: L12C4B22-82-TL
% Mức: NB
Trình bày cách nhận biết và phương pháp giải dạng “Tính toán nhiệt lượng vật thu vào hoặc tỏa ra”. Phân tích nhận định sau và sửa lại nếu sai: “Với cùng m và Δ $T$, chất có c nhỏ hơn cần nhiệt lượng lớn hơn.”
\loigiai{
Trước hết xác định hiện tượng, trạng thái và các đại lượng đã cho. Nêu định nghĩa hoặc hệ thức phù hợp, đổi về đơn vị SI, áp dụng đúng điều kiện và kiểm tra ý nghĩa kết quả. Nhận định cần đối chiếu: Q=mcΔ $T$ khi không có chuyển thể. Vì vậy nhận định đã cho sai; phát biểu đúng là: Q=mcΔ $T$ khi không có chuyển thể.
}
\end{bt}

% ===== Câu 101 =====
% ID: L12C4B22-83-DS
% Mức: NB
\begin{ex}
% ID: L12C4B22-83-DS
% Mức: NB
Xét các phát biểu sau về dạng “Tính toán nhiệt lượng vật thu vào hoặc tỏa ra” (bộ số 4).
\choiceTF
{\True Q=mcΔ $T$ khi không có chuyển thể.}
{Với cùng m và Δ $T$, chất có c nhỏ hơn cần nhiệt lượng lớn hơn.}
{\True Khi giải dạng “Tính toán nhiệt lượng vật thu vào hoặc tỏa ra”, cần xác định đúng trạng thái đầu, trạng thái cuối và quy ước dấu hoặc điều kiện áp dụng.}
{Có thể bỏ qua mọi dữ kiện về khối lượng, nhiệt độ và hiệu suất trong mọi bài toán thuộc dạng này.}
\loigiai{
\begin{itemchoice}
\item (Đúng) Q=mcΔ $T$ khi không có chuyển thể. (Vì): Q=mcΔ $T$ khi không có chuyển thể. b) S: Với cùng m và Δ $T$, chất có c nhỏ hơn cần nhiệt lượng lớn hơn. c) Đ: Khi giải dạng “Tính toán nhiệt lượng vật thu vào hoặc tỏa ra”, cần xác định đúng trạng thái đầu, trạng thái cuối và quy ước dấu hoặc điều kiện áp dụng. d) S: Có thể bỏ qua mọi dữ kiện về khối lượng, nhiệt độ và hiệu suất trong mọi bài toán thuộc dạng này. Kết luận dựa trên định nghĩa, điều kiện áp dụng và quy ước trong SGK KNTT
\item (Sai) Với cùng m và Δ $T$, chất có c nhỏ hơn cần nhiệt lượng lớn hơn.
\item (Đúng) Khi giải dạng “Tính toán nhiệt lượng vật thu vào hoặc tỏa ra”, cần xác định đúng trạng thái đầu, trạng thái cuối và quy ước dấu hoặc điều kiện áp dụng.
\item (Sai) Có thể bỏ qua mọi dữ kiện về khối lượng, nhiệt độ và hiệu suất trong mọi bài toán thuộc dạng này.
\end{itemchoice}
}
\end{ex}

% ===== Câu 102 =====
% ID: L12C4B22-84-TL
% Mức: NB
\begin{bt}
% ID: L12C4B22-84-TL
% Mức: NB
Trình bày cách nhận biết và phương pháp giải dạng “Tính toán nhiệt lượng vật thu vào hoặc tỏa ra”. Phân tích nhận định sau và sửa lại nếu sai: “Với cùng m và Δ $T$, chất có c nhỏ hơn cần nhiệt lượng lớn hơn.”
\loigiai{
Trước hết xác định hiện tượng, trạng thái và các đại lượng đã cho. Nêu định nghĩa hoặc hệ thức phù hợp, đổi về đơn vị SI, áp dụng đúng điều kiện và kiểm tra ý nghĩa kết quả. Nhận định cần đối chiếu: Q=mcΔ $T$ khi không có chuyển thể. Vì vậy nhận định đã cho sai; phát biểu đúng là: Q=mcΔ $T$ khi không có chuyển thể.
}
\end{bt}

% ===== Câu 105 =====
% ID: L12C4B22-85-DS
% Mức: TH
\begin{ex}
% ID: L12C4B22-85-DS
% Mức: TH
Xét các phát biểu sau về dạng “Tính toán nhiệt lượng vật thu vào hoặc tỏa ra” (bộ số 1).
\choiceTF
{Với cùng m và Δ $T$, chất có c nhỏ hơn cần nhiệt lượng lớn hơn.}
{\True Khi giải dạng “Tính toán nhiệt lượng vật thu vào hoặc tỏa ra”, cần xác định đúng trạng thái đầu, trạng thái cuối và quy ước dấu hoặc điều kiện áp dụng.}
{Có thể bỏ qua mọi dữ kiện về khối lượng, nhiệt độ và hiệu suất trong mọi bài toán thuộc dạng này.}
{\True Q=mcΔ $T$ khi không có chuyển thể.}
\loigiai{
\begin{itemchoice}
\item (Sai) Với cùng m và Δ $T$, chất có c nhỏ hơn cần nhiệt lượng lớn hơn. (Vì): Với cùng m và Δ $T$, chất có c nhỏ hơn cần nhiệt lượng lớn hơn. b) Đ: Khi giải dạng “Tính toán nhiệt lượng vật thu vào hoặc tỏa ra”, cần xác định đúng trạng thái đầu, trạng thái cuối và quy ước dấu hoặc điều kiện áp dụng. c) S: Có thể bỏ qua mọi dữ kiện về khối lượng, nhiệt độ và hiệu suất trong mọi bài toán thuộc dạng này. d) Đ: Q=mcΔ $T$ khi không có chuyển thể. Kết luận dựa trên định nghĩa, điều kiện áp dụng và quy ước trong SGK KNTT
\item (Đúng) Khi giải dạng “Tính toán nhiệt lượng vật thu vào hoặc tỏa ra”, cần xác định đúng trạng thái đầu, trạng thái cuối và quy ước dấu hoặc điều kiện áp dụng.
\item (Sai) Có thể bỏ qua mọi dữ kiện về khối lượng, nhiệt độ và hiệu suất trong mọi bài toán thuộc dạng này.
\item (Đúng) Q=mcΔ $T$ khi không có chuyển thể.
\end{itemchoice}
}
\end{ex}

% ===== Câu 106 =====
% ID: L12C4B22-86-TL
% Mức: TH
\begin{bt}
% ID: L12C4B22-86-TL
% Mức: TH
Trình bày cách nhận biết và phương pháp giải dạng “Tính toán nhiệt lượng vật thu vào hoặc tỏa ra”. Phân tích nhận định sau và sửa lại nếu sai: “Q=mcΔ $T$ khi không có chuyển thể.”
\loigiai{
Trước hết xác định hiện tượng, trạng thái và các đại lượng đã cho. Nêu định nghĩa hoặc hệ thức phù hợp, đổi về đơn vị SI, áp dụng đúng điều kiện và kiểm tra ý nghĩa kết quả. Nhận định cần đối chiếu: Q=mcΔ $T$ khi không có chuyển thể. Vì vậy nhận định đã cho đúng.
}
\end{bt}

% ===== Câu 109 =====
% ID: L12C4B22-87-DS
% Mức: TH
\begin{ex}
% ID: L12C4B22-87-DS
% Mức: TH
Xét các phát biểu sau về dạng “Tính toán nhiệt lượng vật thu vào hoặc tỏa ra” (bộ số 5).
\choiceTF
{Với cùng m và Δ $T$, chất có c nhỏ hơn cần nhiệt lượng lớn hơn.}
{\True Khi giải dạng “Tính toán nhiệt lượng vật thu vào hoặc tỏa ra”, cần xác định đúng trạng thái đầu, trạng thái cuối và quy ước dấu hoặc điều kiện áp dụng.}
{Có thể bỏ qua mọi dữ kiện về khối lượng, nhiệt độ và hiệu suất trong mọi bài toán thuộc dạng này.}
{\True Q=mcΔ $T$ khi không có chuyển thể.}
\loigiai{
\begin{itemchoice}
\item (Sai) Với cùng m và Δ $T$, chất có c nhỏ hơn cần nhiệt lượng lớn hơn. (Vì): Với cùng m và Δ $T$, chất có c nhỏ hơn cần nhiệt lượng lớn hơn. b) Đ: Khi giải dạng “Tính toán nhiệt lượng vật thu vào hoặc tỏa ra”, cần xác định đúng trạng thái đầu, trạng thái cuối và quy ước dấu hoặc điều kiện áp dụng. c) S: Có thể bỏ qua mọi dữ kiện về khối lượng, nhiệt độ và hiệu suất trong mọi bài toán thuộc dạng này. d) Đ: Q=mcΔ $T$ khi không có chuyển thể. Kết luận dựa trên định nghĩa, điều kiện áp dụng và quy ước trong SGK KNTT
\item (Đúng) Khi giải dạng “Tính toán nhiệt lượng vật thu vào hoặc tỏa ra”, cần xác định đúng trạng thái đầu, trạng thái cuối và quy ước dấu hoặc điều kiện áp dụng.
\item (Sai) Có thể bỏ qua mọi dữ kiện về khối lượng, nhiệt độ và hiệu suất trong mọi bài toán thuộc dạng này.
\item (Đúng) Q=mcΔ $T$ khi không có chuyển thể.
\end{itemchoice}
}
\end{ex}

% ===== Câu 110 =====
% ID: L12C4B22-88-TL
% Mức: VD
\begin{bt}
% ID: L12C4B22-88-TL
% Mức: VD
Trình bày cách nhận biết và phương pháp giải dạng “Tính toán nhiệt lượng vật thu vào hoặc tỏa ra”. Phân tích nhận định sau và sửa lại nếu sai: “Với cùng m và Δ $T$, chất có c nhỏ hơn cần nhiệt lượng lớn hơn.”
\loigiai{
Trước hết xác định hiện tượng, trạng thái và các đại lượng đã cho. Nêu định nghĩa hoặc hệ thức phù hợp, đổi về đơn vị SI, áp dụng đúng điều kiện và kiểm tra ý nghĩa kết quả. Nhận định cần đối chiếu: Q=mcΔ $T$ khi không có chuyển thể. Vì vậy nhận định đã cho sai; phát biểu đúng là: Q=mcΔ $T$ khi không có chuyển thể.
}
\end{bt}

% ===== Câu 113 =====
% ID: L12C4B22-89-DS
% Mức: VD
\begin{ex}
% ID: L12C4B22-89-DS
% Mức: VD
Xét các phát biểu sau về dạng “Tính toán nhiệt lượng vật thu vào hoặc tỏa ra” (bộ số 2).
\choiceTF
{\True Khi giải dạng “Tính toán nhiệt lượng vật thu vào hoặc tỏa ra”, cần xác định đúng trạng thái đầu, trạng thái cuối và quy ước dấu hoặc điều kiện áp dụng.}
{Có thể bỏ qua mọi dữ kiện về khối lượng, nhiệt độ và hiệu suất trong mọi bài toán thuộc dạng này.}
{\True Q=mcΔ $T$ khi không có chuyển thể.}
{Với cùng m và Δ $T$, chất có c nhỏ hơn cần nhiệt lượng lớn hơn.}
\loigiai{
\begin{itemchoice}
\item (Đúng) Khi giải dạng “Tính toán nhiệt lượng vật thu vào hoặc tỏa ra”, cần xác định đúng trạng thái đầu, trạng thái cuối và quy ước dấu hoặc điều kiện áp dụng. (Vì): Khi giải dạng “Tính toán nhiệt lượng vật thu vào hoặc tỏa ra”, cần xác định đúng trạng thái đầu, trạng thái cuối và quy ước dấu hoặc điều kiện áp dụng. b) S: Có thể bỏ qua mọi dữ kiện về khối lượng, nhiệt độ và hiệu suất trong mọi bài toán thuộc dạng này. c) Đ: Q=mcΔ $T$ khi không có chuyển thể. d) S: Với cùng m và Δ $T$, chất có c nhỏ hơn cần nhiệt lượng lớn hơn. Kết luận dựa trên định nghĩa, điều kiện áp dụng và quy ước trong SGK KNTT
\item (Sai) Có thể bỏ qua mọi dữ kiện về khối lượng, nhiệt độ và hiệu suất trong mọi bài toán thuộc dạng này.
\item (Đúng) Q=mcΔ $T$ khi không có chuyển thể.
\item (Sai) Với cùng m và Δ $T$, chất có c nhỏ hơn cần nhiệt lượng lớn hơn.
\end{itemchoice}
}
\end{ex}

% ===== Câu 114 =====
% ID: L12C4B22-90-TL
% Mức: VDC
\begin{bt}
% ID: L12C4B22-90-TL
% Mức: VDC
Trình bày cách nhận biết và phương pháp giải dạng “Tính toán nhiệt lượng vật thu vào hoặc tỏa ra”. Phân tích nhận định sau và sửa lại nếu sai: “Q=mcΔ $T$ khi không có chuyển thể.”
\loigiai{
Trước hết xác định hiện tượng, trạng thái và các đại lượng đã cho. Nêu định nghĩa hoặc hệ thức phù hợp, đổi về đơn vị SI, áp dụng đúng điều kiện và kiểm tra ý nghĩa kết quả. Nhận định cần đối chiếu: Q=mcΔ $T$ khi không có chuyển thể. Vì vậy nhận định đã cho đúng.
}
\end{bt}

% ===== Câu 117 =====
% ID: L12C4B22-91-DS
% Mức: VDC
\begin{ex}
% ID: L12C4B22-91-DS
% Mức: VDC
Xét các phát biểu sau về dạng “Tính toán nhiệt lượng vật thu vào hoặc tỏa ra” (bộ số 3).
\choiceTF
{Có thể bỏ qua mọi dữ kiện về khối lượng, nhiệt độ và hiệu suất trong mọi bài toán thuộc dạng này.}
{\True Q=mcΔ $T$ khi không có chuyển thể.}
{Với cùng m và Δ $T$, chất có c nhỏ hơn cần nhiệt lượng lớn hơn.}
{\True Khi giải dạng “Tính toán nhiệt lượng vật thu vào hoặc tỏa ra”, cần xác định đúng trạng thái đầu, trạng thái cuối và quy ước dấu hoặc điều kiện áp dụng.}
\loigiai{
\begin{itemchoice}
\item (Sai) Có thể bỏ qua mọi dữ kiện về khối lượng, nhiệt độ và hiệu suất trong mọi bài toán thuộc dạng này. (Vì): Có thể bỏ qua mọi dữ kiện về khối lượng, nhiệt độ và hiệu suất trong mọi bài toán thuộc dạng này. b) Đ: Q=mcΔ $T$ khi không có chuyển thể. c) S: Với cùng m và Δ $T$, chất có c nhỏ hơn cần nhiệt lượng lớn hơn. d) Đ: Khi giải dạng “Tính toán nhiệt lượng vật thu vào hoặc tỏa ra”, cần xác định đúng trạng thái đầu, trạng thái cuối và quy ước dấu hoặc điều kiện áp dụng. Kết luận dựa trên định nghĩa, điều kiện áp dụng và quy ước trong SGK KNTT
\item (Đúng) Q=mcΔ $T$ khi không có chuyển thể.
\item (Sai) Với cùng m và Δ $T$, chất có c nhỏ hơn cần nhiệt lượng lớn hơn.
\item (Đúng) Khi giải dạng “Tính toán nhiệt lượng vật thu vào hoặc tỏa ra”, cần xác định đúng trạng thái đầu, trạng thái cuối và quy ước dấu hoặc điều kiện áp dụng.
\end{itemchoice}
}
\end{ex}

% ===== Câu 118 =====
% ID: L12C4B22-92-TL
% Mức: VDC
\begin{bt}
% ID: L12C4B22-92-TL
% Mức: VDC
Trình bày cách nhận biết và phương pháp giải dạng “Tính toán nhiệt lượng vật thu vào hoặc tỏa ra”. Phân tích nhận định sau và sửa lại nếu sai: “Q=mcΔ $T$ khi không có chuyển thể.”
\loigiai{
Trước hết xác định hiện tượng, trạng thái và các đại lượng đã cho. Nêu định nghĩa hoặc hệ thức phù hợp, đổi về đơn vị SI, áp dụng đúng điều kiện và kiểm tra ý nghĩa kết quả. Nhận định cần đối chiếu: Q=mcΔ $T$ khi không có chuyển thể. Vì vậy nhận định đã cho đúng.
}
\end{bt}

% ===== Câu 123 =====
% ID: L12C4B22-93-TLN
% Mức: NB
\begin{ex}
% ID: L12C4B22-93-TLN
% Mức: NB
Cho bốn nhận định: (1) Q=mcΔ $T$ khi không có chuyển thể. (2) Với cùng m và Δ $T$, chất có c nhỏ hơn cần nhiệt lượng lớn hơn. (3) Cần kiểm tra đơn vị và điều kiện áp dụng khi xử lí dạng “Tính toán nhiệt lượng vật thu vào hoặc tỏa ra”. (4) Có thể thay tùy ý nhiệt độ Celsius cho nhiệt độ tuyệt đối trong mọi công thức. Có bao nhiêu nhận định đúng? Trả lời bằng một số nguyên.
\shortans{2}
\loigiai{
Nhận định (1) và (3) đúng; (2) và (4) sai. Vậy có 2 nhận định đúng.
}
\end{ex}

% ===== Câu 127 =====
% ID: L12C4B22-94-TLN
% Mức: TH
\begin{ex}
% ID: L12C4B22-94-TLN
% Mức: TH
Cho bốn nhận định: (1) Q=mcΔ $T$ khi không có chuyển thể. (2) Với cùng m và Δ $T$, chất có c nhỏ hơn cần nhiệt lượng lớn hơn. (3) Cần kiểm tra đơn vị và điều kiện áp dụng khi xử lí dạng “Tính toán nhiệt lượng vật thu vào hoặc tỏa ra”. (4) Có thể thay tùy ý nhiệt độ Celsius cho nhiệt độ tuyệt đối trong mọi công thức. Có bao nhiêu nhận định đúng? Trả lời bằng một số nguyên.
\shortans{2}
\loigiai{
Nhận định (1) và (3) đúng; (2) và (4) sai. Vậy có 2 nhận định đúng.
}
\end{ex}

% ===== Câu 131 =====
% ID: L12C4B22-95-TLN
% Mức: VD
\begin{ex}
% ID: L12C4B22-95-TLN
% Mức: VD
Cho bốn nhận định: (1) Q=mcΔ $T$ khi không có chuyển thể. (2) Với cùng m và Δ $T$, chất có c nhỏ hơn cần nhiệt lượng lớn hơn. (3) Cần kiểm tra đơn vị và điều kiện áp dụng khi xử lí dạng “Tính toán nhiệt lượng vật thu vào hoặc tỏa ra”. (4) Có thể thay tùy ý nhiệt độ Celsius cho nhiệt độ tuyệt đối trong mọi công thức. Có bao nhiêu nhận định đúng? Trả lời bằng một số nguyên.
\shortans{2}
\loigiai{
Nhận định (1) và (3) đúng; (2) và (4) sai. Vậy có 2 nhận định đúng.
}
\end{ex}

% ===== Câu 135 =====
% ID: L12C4B22-96-TLN
% Mức: VD
\begin{ex}
% ID: L12C4B22-96-TLN
% Mức: VD
Cho bốn nhận định: (1) Q=mcΔ $T$ khi không có chuyển thể. (2) Với cùng m và Δ $T$, chất có c nhỏ hơn cần nhiệt lượng lớn hơn. (3) Cần kiểm tra đơn vị và điều kiện áp dụng khi xử lí dạng “Tính toán nhiệt lượng vật thu vào hoặc tỏa ra”. (4) Có thể thay tùy ý nhiệt độ Celsius cho nhiệt độ tuyệt đối trong mọi công thức. Có bao nhiêu nhận định đúng? Trả lời bằng một số nguyên.
\shortans{2}
\loigiai{
Nhận định (1) và (3) đúng; (2) và (4) sai. Vậy có 2 nhận định đúng.
}
\end{ex}

% ===== Câu 139 =====
% ID: L12C4B22-97-TLN
% Mức: VDC
\begin{ex}
% ID: L12C4B22-97-TLN
% Mức: VDC
Cho bốn nhận định: (1) Q=mcΔ $T$ khi không có chuyển thể. (2) Với cùng m và Δ $T$, chất có c nhỏ hơn cần nhiệt lượng lớn hơn. (3) Cần kiểm tra đơn vị và điều kiện áp dụng khi xử lí dạng “Tính toán nhiệt lượng vật thu vào hoặc tỏa ra”. (4) Có thể thay tùy ý nhiệt độ Celsius cho nhiệt độ tuyệt đối trong mọi công thức. Có bao nhiêu nhận định đúng? Trả lời bằng một số nguyên.
\shortans{2}
\loigiai{
Nhận định (1) và (3) đúng; (2) và (4) sai. Vậy có 2 nhận định đúng.
}
\end{ex}

% ===== Câu 142 =====
% ID: L12C4B22-98-TLN
% Mức: VDC
\begin{ex}
% ID: L12C4B22-98-TLN
% Mức: VDC
Cho bốn nhận định: (1) Q=mcΔ $T$ khi không có chuyển thể. (2) Với cùng m và Δ $T$, chất có c nhỏ hơn cần nhiệt lượng lớn hơn. (3) Cần kiểm tra đơn vị và điều kiện áp dụng khi xử lí dạng “Tính toán nhiệt lượng vật thu vào hoặc tỏa ra”. (4) Có thể thay tùy ý nhiệt độ Celsius cho nhiệt độ tuyệt đối trong mọi công thức. Có bao nhiêu nhận định đúng? Trả lời bằng một số nguyên.
\shortans{2}
\loigiai{
Nhận định (1) và (3) đúng; (2) và (4) sai. Vậy có 2 nhận định đúng.
}
\end{ex}

% ===== Câu 166 =====
% ID: L12C4B22-99-TN
% Mức: NB
\begin{ex}
% ID: L12C4B22-99-TN
% Mức: NB
Câu 1 thuộc dạng “Tính toán nhiệt lượng vật thu vào hoặc tỏa ra”. Phát biểu nào sau đây đúng?
\choice
{Với cùng m và Δ $T$, chất có c nhỏ hơn cần nhiệt lượng lớn hơn.}
{Nhiệt lượng và nhiệt độ là hai đại lượng hoàn toàn giống nhau.}
{Mọi quá trình nhiệt học đều không phụ thuộc điều kiện thực hiện.}
{\True Q=mcΔ $T$ khi không có chuyển thể.}
\loigiai{
Phát biểu đúng là: Q=mcΔ $T$ khi không có chuyển thể. Các phương án còn lại trái với mô hình, định nghĩa hoặc hệ thức nhiệt học tương ứng.
}
\end{ex}

% ===== Câu 167 =====
% ID: L12C4B22-100-TN
% Mức: NB
\begin{ex}
% ID: L12C4B22-100-TN
% Mức: NB
Câu 5 thuộc dạng “Tính toán nhiệt lượng vật thu vào hoặc tỏa ra”. Phát biểu nào sau đây đúng?
\choice
{Với cùng m và Δ $T$, chất có c nhỏ hơn cần nhiệt lượng lớn hơn.}
{Nhiệt lượng và nhiệt độ là hai đại lượng hoàn toàn giống nhau.}
{Mọi quá trình nhiệt học đều không phụ thuộc điều kiện thực hiện.}
{\True Q=mcΔ $T$ khi không có chuyển thể.}
\loigiai{
Phát biểu đúng là: Q=mcΔ $T$ khi không có chuyển thể. Các phương án còn lại trái với mô hình, định nghĩa hoặc hệ thức nhiệt học tương ứng.
}
\end{ex}

% ===== Câu 168 =====
% ID: L12C4B22-101-TN
% Mức: NB
\begin{ex}
% ID: L12C4B22-101-TN
% Mức: NB
Câu 9 thuộc dạng “Tính toán nhiệt lượng vật thu vào hoặc tỏa ra”. Phát biểu nào sau đây đúng?
\choice
{Với cùng m và Δ $T$, chất có c nhỏ hơn cần nhiệt lượng lớn hơn.}
{Nhiệt lượng và nhiệt độ là hai đại lượng hoàn toàn giống nhau.}
{Mọi quá trình nhiệt học đều không phụ thuộc điều kiện thực hiện.}
{\True Q=mcΔ $T$ khi không có chuyển thể.}
\loigiai{
Phát biểu đúng là: Q=mcΔ $T$ khi không có chuyển thể. Các phương án còn lại trái với mô hình, định nghĩa hoặc hệ thức nhiệt học tương ứng.
}
\end{ex}

% ===== Câu 169 =====
% ID: L12C4B22-102-TN
% Mức: TH
\begin{ex}
% ID: L12C4B22-102-TN
% Mức: TH
Câu 2 thuộc dạng “Tính toán nhiệt lượng vật thu vào hoặc tỏa ra”. Phát biểu nào sau đây đúng?
\choice
{Nhiệt lượng và nhiệt độ là hai đại lượng hoàn toàn giống nhau.}
{Mọi quá trình nhiệt học đều không phụ thuộc điều kiện thực hiện.}
{\True Q=mcΔ $T$ khi không có chuyển thể.}
{Với cùng m và Δ $T$, chất có c nhỏ hơn cần nhiệt lượng lớn hơn.}
\loigiai{
Phát biểu đúng là: Q=mcΔ $T$ khi không có chuyển thể. Các phương án còn lại trái với mô hình, định nghĩa hoặc hệ thức nhiệt học tương ứng.
}
\end{ex}

% ===== Câu 170 =====
% ID: L12C4B22-103-TN
% Mức: TH
\begin{ex}
% ID: L12C4B22-103-TN
% Mức: TH
Câu 6 thuộc dạng “Tính toán nhiệt lượng vật thu vào hoặc tỏa ra”. Phát biểu nào sau đây đúng?
\choice
{Nhiệt lượng và nhiệt độ là hai đại lượng hoàn toàn giống nhau.}
{Mọi quá trình nhiệt học đều không phụ thuộc điều kiện thực hiện.}
{\True Q=mcΔ $T$ khi không có chuyển thể.}
{Với cùng m và Δ $T$, chất có c nhỏ hơn cần nhiệt lượng lớn hơn.}
\loigiai{
Phát biểu đúng là: Q=mcΔ $T$ khi không có chuyển thể. Các phương án còn lại trái với mô hình, định nghĩa hoặc hệ thức nhiệt học tương ứng.
}
\end{ex}

% ===== Câu 171 =====
% ID: L12C4B22-104-TN
% Mức: TH
\begin{ex}
% ID: L12C4B22-104-TN
% Mức: TH
Câu 10 thuộc dạng “Tính toán nhiệt lượng vật thu vào hoặc tỏa ra”. Phát biểu nào sau đây đúng?
\choice
{Nhiệt lượng và nhiệt độ là hai đại lượng hoàn toàn giống nhau.}
{Mọi quá trình nhiệt học đều không phụ thuộc điều kiện thực hiện.}
{\True Q=mcΔ $T$ khi không có chuyển thể.}
{Với cùng m và Δ $T$, chất có c nhỏ hơn cần nhiệt lượng lớn hơn.}
\loigiai{
Phát biểu đúng là: Q=mcΔ $T$ khi không có chuyển thể. Các phương án còn lại trái với mô hình, định nghĩa hoặc hệ thức nhiệt học tương ứng.
}
\end{ex}

% ===== Câu 172 =====
% ID: L12C4B22-105-TN
% Mức: VD
\begin{ex}
% ID: L12C4B22-105-TN
% Mức: VD
Câu 3 thuộc dạng “Tính toán nhiệt lượng vật thu vào hoặc tỏa ra”. Phát biểu nào sau đây đúng?
\choice
{Mọi quá trình nhiệt học đều không phụ thuộc điều kiện thực hiện.}
{\True Q=mcΔ $T$ khi không có chuyển thể.}
{Với cùng m và Δ $T$, chất có c nhỏ hơn cần nhiệt lượng lớn hơn.}
{Nhiệt lượng và nhiệt độ là hai đại lượng hoàn toàn giống nhau.}
\loigiai{
Phát biểu đúng là: Q=mcΔ $T$ khi không có chuyển thể. Các phương án còn lại trái với mô hình, định nghĩa hoặc hệ thức nhiệt học tương ứng.
}
\end{ex}

% ===== Câu 173 =====
% ID: L12C4B22-106-TN
% Mức: VD
\begin{ex}
% ID: L12C4B22-106-TN
% Mức: VD
Câu 7 thuộc dạng “Tính toán nhiệt lượng vật thu vào hoặc tỏa ra”. Phát biểu nào sau đây đúng?
\choice
{Mọi quá trình nhiệt học đều không phụ thuộc điều kiện thực hiện.}
{\True Q=mcΔ $T$ khi không có chuyển thể.}
{Với cùng m và Δ $T$, chất có c nhỏ hơn cần nhiệt lượng lớn hơn.}
{Nhiệt lượng và nhiệt độ là hai đại lượng hoàn toàn giống nhau.}
\loigiai{
Phát biểu đúng là: Q=mcΔ $T$ khi không có chuyển thể. Các phương án còn lại trái với mô hình, định nghĩa hoặc hệ thức nhiệt học tương ứng.
}
\end{ex}

% ===== Câu 174 =====
% ID: L12C4B22-107-TN
% Mức: VDC
\begin{ex}
% ID: L12C4B22-107-TN
% Mức: VDC
Câu 4 thuộc dạng “Tính toán nhiệt lượng vật thu vào hoặc tỏa ra”. Phát biểu nào sau đây đúng?
\choice
{\True Q=mcΔ $T$ khi không có chuyển thể.}
{Với cùng m và Δ $T$, chất có c nhỏ hơn cần nhiệt lượng lớn hơn.}
{Nhiệt lượng và nhiệt độ là hai đại lượng hoàn toàn giống nhau.}
{Mọi quá trình nhiệt học đều không phụ thuộc điều kiện thực hiện.}
\loigiai{
Phát biểu đúng là: Q=mcΔ $T$ khi không có chuyển thể. Các phương án còn lại trái với mô hình, định nghĩa hoặc hệ thức nhiệt học tương ứng.
}
\end{ex}

% ===== Câu 175 =====
% ID: L12C4B22-108-TN
% Mức: VDC
\begin{ex}
% ID: L12C4B22-108-TN
% Mức: VDC
Câu 8 thuộc dạng “Tính toán nhiệt lượng vật thu vào hoặc tỏa ra”. Phát biểu nào sau đây đúng?
\choice
{\True Q=mcΔ $T$ khi không có chuyển thể.}
{Với cùng m và Δ $T$, chất có c nhỏ hơn cần nhiệt lượng lớn hơn.}
{Nhiệt lượng và nhiệt độ là hai đại lượng hoàn toàn giống nhau.}
{Mọi quá trình nhiệt học đều không phụ thuộc điều kiện thực hiện.}
\loigiai{
Phát biểu đúng là: Q=mcΔ $T$ khi không có chuyển thể. Các phương án còn lại trái với mô hình, định nghĩa hoặc hệ thức nhiệt học tương ứng.
}
\end{ex}
